\chapterhead{6}{Finding the Missing Arrow}{Using leftover arrows to explain what a thing does}\label{ch6}

A force diagram is a model. It keeps the direction and the relative size of each force, and it drops the exact shape of the thing, its colour and its parts. It is good when the arrows explain what the thing did. If they do not, an arrow is missing.

\begin{center}
\sffamily\small\renewcommand{\arraystretch}{1.5}
\begin{tabular}{@{}p{3.4cm}p{3.7cm}p{3.7cm}@{}}
\toprule
\textbf{Model} & \textbf{What it keeps} & \textbf{What it drops}\\
\midrule
An arrow for a force & direction and size & the exact shape of the thing\\
Popping the arrows & the leftover push & which arrow came from where\\
Friction arrow & points against sliding & how rough each bump is\\
Coins times distance & when a seesaw balances & the weight of the ruler\\
A ramp picture & how steep and how long & the bumps on the board\\
\bottomrule
\end{tabular}
\end{center}

\subhead{Three rules for reading a diagram}
\begin{enumerate}[leftmargin=1.6em,itemsep=0.2em]
\item Arrows that are equal and opposite pop to nothing. The thing keeps doing what it was doing: staying still, or moving at a steady speed.
\item A leftover arrow means the thing speeds up, slows down or turns, in the direction of the leftover.
\item If the thing does something the arrows cannot explain, look for the missing arrow.
\end{enumerate}

\begin{worked}{A crate on a floor (practice numbers)}
You push a crate right with an arrow of size 8. Friction pulls left with size 5. Pop them: $8-5=3$ to the right. The crate speeds up, slowly.\\
Now push with 5. Pop: $5-5=0$. The crate slides at a steady speed, or stays put.
\end{worked}

\begin{center}
\begin{tikzpicture}
  \node[bigcir,minimum size=2.2cm] (a) at (0,0) {}; \node[font=\sffamily\scriptsize,text width=1.9cm,align=center] at (0,0) {push 8\\friction 5};
  \node[valcir] (b) at (5,0) {3 right};
  \draw[popA] (a) -- node[above,note] {pop} (b);
  \node[note,align=left,anchor=west,text width=3.6cm] at (6.1,0) {\textbf{kept:} which way and how much\\\textbf{lost:} which arrow is which};
\end{tikzpicture}
\end{center}

\begin{multix}[hmm]\multixbot[hmm]\tcblower
A sled glides along flat ice at a steady speed, so it must have a leftover push forward to keep it going. Does it?
\end{multix}

\trust{We checked each diagram by asking whether its arrows explain what the thing did.}
{When the arrows popped to the same result we saw, the diagram explained it. When they did not, we found the arrow that was missing.}
{That every force is shown. A diagram can pass on one test and still leave out a small force, such as air.}

\begin{keepbox}
We turned a sliding crate into two arrows. \textbf{We kept} \wline[3.4cm] \textbf{and we set aside} \wline[3.4cm].
\end{keepbox}

\clearpage

\begin{puzzle}{Spot the missing arrow}
A box slides across a floor at a steady speed. Only one arrow is drawn: push right, size 4.
\begin{center}
\begin{tikzpicture}
  \fill[sky!50,draw=ink,line width=1pt] (0,0) rectangle (1.4,0.9);
  \draw[->,line width=1.8pt,color=leaf!70!black] (1.4,0.45) -- (3.0,0.45) node[right,font=\sffamily\scriptsize] {push 4};
  \draw[line width=1pt] (-1,0) -- (4.5,0);
\end{tikzpicture}
\end{center}
Which arrow is missing, which way does it point, and how big is it? \wline[4cm]
\end{puzzle}

\begin{puzzle}{A ball on a table}
A ball sits still on a table. Only gravity is drawn, pointing down, size 3. What is missing? \wline[3cm]\\
How big is it? \wline[1.5cm]
\end{puzzle}

\begin{puzzle}{Pop them}
Pop these arrows and write the leftover.\\
(a) 7 right, 3 left, 2 left: \wline[2.4cm]\\
(b) 5 up, 5 down: \wline[2.4cm]\\
(c) 4 right, 9 left: \wline[2.4cm]
\end{puzzle}

\begin{puzzle}{Tug of war}
Three children pull a rope left, each with size 4. Two children pull it right, each with size 5. Which way does the rope go, and with what leftover? \wline[4cm]
\end{puzzle}

\begin{puzzle}{The floating boat}
A toy boat floats still. Gravity pulls down with size 10. How big is the water's push, and which way? \wline[3.5cm]
\end{puzzle}

\clearpage

\begin{puzzle}{Choose the surface}
From Chapter 2, a block slid on carpet 19~cm, on card 46~cm and on a plastic tray 78~cm. You want a sled to stop soonest. Which surface? \wline[2.5cm]\\
Why does it stop soonest there? \wline[3cm]
\end{puzzle}

\begin{puzzle}{A seesaw mystery}
A ruler seesaw balances with 4 coins at 5~cm on the left and a bag on the right at 10~cm. How many coins does the bag weigh? \wline[1.8cm]
\end{puzzle}

\begin{puzzle}{Skater}
A skater glides at a steady speed on flat ice. How many forward arrows are needed to explain this? \wline[2cm]\\
Name the two arrows that do pop to nothing, one up and one down. \wline[4cm]
\end{puzzle}

\begin{puzzle}{Arrows for a leaf}
A leaf drifts down at a steady speed. Gravity pulls down with size 2. How big is the air's push, and which way? \wline[3.5cm]
\end{puzzle}

\begin{puzzle}{Pop them again}
9 right, 4 left, 5 left. Leftover: \wline[2.4cm]
\end{puzzle}

\begin{puzzle}{The speeding car}
A toy car rolls down a ramp. Gravity pulls down and the ramp pushes out. A small leftover points down the slope. Does the car speed up, slow down or stay steady? \wline[3cm]
\end{puzzle}

\begin{puzzle}{Another seesaw}
Left: 2 coins at 12~cm. Right: 4 coins at 6~cm. Does it balance? \wline[2cm] Show the numbers. \wline[3cm]
\end{puzzle}

\begin{puzzle}{A neat diagram}
Ben says his force diagram must be right because it looks neat. What should he check instead? \wline[5cm]
\end{puzzle}

\clearpage
\begin{puzzle}{Draw the arrows: a parachute}
A toy parachute drifts down at a steady speed. Draw the arrows from the dot and label each size.
\begin{center}\begin{tikzpicture}\draw[line width=1pt,rounded corners=3pt] (0,0) rectangle (5,3); \fill[ink] (2.5,1.5) circle (3pt);\end{tikzpicture}\end{center}
\end{puzzle}

\begin{puzzle}{Draw the arrows: a speeding box}
A box is pushed right with size 7 across a floor. Friction is size 3. Draw the horizontal arrows and write the leftover. \wline[3cm]
\begin{center}\begin{tikzpicture}\draw[line width=1pt,rounded corners=3pt] (0,0) rectangle (5,3); \fill[ink] (2.5,1.5) circle (3pt);\end{tikzpicture}\end{center}
\end{puzzle}

\begin{puzzle}{Draw the arrows: a floating boat}
A boat floats still on a pond. Draw the arrows from the dot and say why they pop to nothing. \wline[4cm]
\begin{center}\begin{tikzpicture}\draw[line width=1pt,rounded corners=3pt] (0,0) rectangle (5,3); \fill[ink] (2.5,1.5) circle (3pt);\end{tikzpicture}\end{center}
\end{puzzle}

\clearpage
\clearpage
\thispagestyle{fancy}\label{slop}
\subhead{Slop check}
Slop is hasty, gaudy work that forgot to take care. It can be made by a person or a machine, quickly or slowly. What matters is whether the work shows care for three things:
\begin{itemize}[leftmargin=1.6em]
\item \textbf{Salient:} Is the main point easy to see?
\item \textbf{Simple:} Is there anything that does nothing but get in the way?
\item \textbf{Relevant:} Does everything on the page help answer the question?
\end{itemize}

Here is a hasty poster about sliding on different surfaces. Study it, then decide what is wrong.

\begin{center}
\begin{tikzpicture}[scale=0.97]
  \fill[yellow!40,draw=purple!80!black,line width=3pt] (0,0) rectangle (11.4,7.6);
  \node[star,star points=12,fill=red!60,draw=blue,line width=2pt,minimum size=2.2cm,font=\bfseries\scriptsize,rotate=10,text=yellow] at (1.4,6.2) {WOW!!!};
  \node[font=\itshape\bfseries\Large,text=teal,rotate=-4] at (6.4,6.8) {FrIcTiOn Is MaGiC};
  \node[font=\ttfamily\tiny,text=orange] at (6.4,6.15) {(blocks maybe? who knows)};
  \foreach \x/\h/\c in {3/1.2/red,5.3/3.2/blue,7.6/5.0/green}{
    \fill[\c,draw=black] (\x,0.9) rectangle (\x+1.3,0.9+\h);
  }
  \node[font=\tiny,text=black,rotate=60] at (9.8,3) {number number number};
  \foreach \x/\y in {1.2/1.5,1.6/4.2,10.2/5.6,0.8/3}{\node[star,star points=5,fill=yellow,draw=red,minimum size=0.6cm] at (\x,\y) {};}
  \node[font=\tiny,text=lime!50!black] at (5.4,0.45) {slide it!!!  best surface ever!!!  wow  wow};
\end{tikzpicture}
\end{center}

\begin{slopbox}
\textbf{(a)} What is the poster's main point? Can you tell?\\[3pt]\rule{\linewidth}{0.5pt}\\[6pt]
\textbf{(b)} Name two things that are not simple, because they only distract.\\[3pt]\rule{\linewidth}{0.5pt}\\[6pt]
\textbf{(c)} Name two things that are missing, so the bars cannot be read.\\[3pt]\rule{\linewidth}{0.5pt}\\[6pt]
\textbf{(d)} Which part is not relevant to sliding at all?\\[3pt]\rule{\linewidth}{0.5pt}
\end{slopbox}

\clearpage
\thispagestyle{fancy}
\subhead{Now you try}
On a separate sheet, redraw a clear bar chart of how far the block slid: carpet 19~cm, card 46~cm and plastic tray 78~cm. Give it a title, label the axes with units, use one or two calm colours, and delete everything that does not help. Check it against the three words: salient, simple, relevant.

\begin{center}
\begin{tikzpicture}
  \draw[ink!70,line width=1pt,rounded corners=3pt] (0,0) rectangle (11.4,1.0);
  \node[font=\sffamily\small,anchor=west] at (0.2,0.5) {Title:};
  \draw[line width=1.3pt] (1,1.6) -- (1,9.6);\draw[line width=1.3pt] (1,1.6) -- (11,1.6);
  \foreach \y in {2.6,3.6,...,8.6}{\draw[gray!40] (1,\y) -- (11,\y);}
  \node[font=\sffamily\scriptsize,rotate=90] at (0.3,5.6) {label and unit here};
  \node[font=\sffamily\scriptsize] at (6,1.3) {label what the axis stands for};
\end{tikzpicture}
\end{center}

\clearpage
\thispagestyle{fancy}
\subhead{What did we keep, across the whole book?}
For each chapter, say one thing that was simplified, what was kept, and what was set aside. A sample is done for you.

\renewcommand{\arraystretch}{2.3}
\begin{center}\sffamily\scriptsize
\begin{tabularx}{\textwidth}{|>{\bfseries}p{2.1cm}|X|X|X|}
\hline
\textbf{Chapter} & \textbf{What we simplified} & \textbf{What we kept} & \textbf{What we set aside} \\ \hline
1 Pushes and pulls & a push drawn as an arrow & direction and size & the shape of the thing \\ \hline
2 Sliding & & & \\ \hline
3 Falling and floating & & & \\ \hline
4 Balance & & & \\ \hline
5 Ramps & & & \\ \hline
6 Missing arrow & & & \\ \hline
\end{tabularx}
\end{center}

\begin{multix}[wow]\multixbot[wow]\tcblower
Everything you drew in this book was a simplification of something. A model is allowed to drop things, as long as you say which. Next time something stops, you will know to look for the arrow.
\end{multix}

\clearpage
\thispagestyle{fancy}
\subhead{Write a force report}
\begin{tcolorbox}[enhanced,colback=white,colframe=ink!60,boxrule=1pt,arc=2mm,left=8pt,right=8pt,top=5pt,bottom=5pt]
\sffamily\small
\begin{enumerate}[leftmargin=1.6em,itemsep=0.2em]
\item \textbf{What I looked at.}
\item \textbf{What I changed, measured and kept the same.}
\item \textbf{My readings.} All of them.
\item \textbf{My answer.} With averages and spread.
\item \textbf{What I kept and what it hides.}
\end{enumerate}
\end{tcolorbox}
\vspace{4pt}
Pick one test from this book and write a report on it. Use the lines below.\par\medskip
\foreach \i in {1,...,12}{\rule{\linewidth}{0.4pt}\\[18pt]}
\clearpage

\thispagestyle{fancy}
\subhead{Tests to try}
Pick one and design it with the habits from this book: change one thing, keep the rest the same, repeat, and average.
\begin{enumerate}[leftmargin=1.6em,itemsep=0.5em]
\item Does the weight of a block change how far it slides down a ramp?
\item Does a wider or longer sheet of paper fall more slowly?
\item Does the length of the pull string matter for the rubber-band stretch?
\item What changes how many coins a clay boat can hold?
\item Does a heavier car roll farther down the same ramp?
\item Does the pivot's position change where a ruler balances?
\end{enumerate}
\medskip
Write the test you pick:\par\medskip
\foreach \i in {1,...,6}{\rule{\linewidth}{0.4pt}\\[16pt]}
Change: \wline[3cm] Measure: \wline[3cm] Keep the same: \wline[3cm]
\clearpage

\thispagestyle{fancy}
\subhead{A week of forces}
Each day, notice one push or pull and draw its arrows.\par\medskip
\renewcommand{\arraystretch}{3.2}
\begin{tabularx}{\linewidth}{|>{\sffamily\small\bfseries}p{1.8cm}|X|X|}
\hline
Day & What pushed or pulled & Arrows I would draw\\ \hline
Mon & & \\ \hline
Tue & & \\ \hline
Wed & & \\ \hline
Thu & & \\ \hline
Fri & & \\ \hline
Sat & & \\ \hline
Sun & & \\ \hline
\end{tabularx}
\clearpage

\thispagestyle{fancy}
\subhead{Questions I still have}
Write what you still wonder about forces. Then pick one and think how to test it.\par\medskip
\foreach \i in {1,...,16}{\rule{\linewidth}{0.4pt}\\[16pt]}
\clearpage
